\section{Inexact delivery has an exact defect certificate}
\label{sec:inexact}

The empirical ``exact rung'' performs repeated $\omega=1$ row operations; it
does not call a block solver.  The following identity states exactly how
close such a phase is to mathematical settlement.

\begin{theorem}[Residual-to-settlement defect identity]
\label{thm:inexact-settlement}
Let $y$ be any point and $\bar y:=\mathcal E_U^\ell(y)$.  Put
$e_U:=e_\ell(y)_U$.  Then
\begin{align}
 \bar y_U-y_U&=Q_{UU}^{-1}e_U,
 \label{eq:inexact-point-defect}\\
 \Psi_\ell(y)-\Psi_\ell(\bar y)
 &=\frac12e_U^TQ_{UU}^{-1}e_U
 \leq\frac1{2\alpha}\|e_U\|_2^2.
 \label{eq:inexact-objective-defect}
\end{align}
If
$|e_u|/\sqrt{d_u}\leq\eta$ for every $u\in U$, then
\begin{equation}
 \Psi_\ell(y)-\Psi_\ell(\bar y)
 \leq\frac{\eta^2\vol(U)}{2\alpha}.
 \label{eq:degree-gate-settlement-defect}
\end{equation}
\end{theorem}

\begin{proof}
The first two displays are
\cref{eq:dirichlet-settlement,eq:settlement-energy-drop}.  Since
$Q_{UU}\succeq\alpha I$,
$Q_{UU}^{-1}\preceq\alpha^{-1}I$.  The coordinatewise hypothesis gives
$\|e_U\|_2^2\leq\eta^2\sum_{u\in U}d_u$.
\end{proof}

Thus a delivery rung supplies two independent objects: a terminal residual
certificate for the verifier and a quantitative settlement defect inside its
touched region.  The first does not automatically imply the second with a
small graph-independent total error unless the touched volume is also
controlled; \cref{eq:degree-gate-settlement-defect} exposes that dependence.
